\subsection{Caché LRU de bloques descomprimidos}

Para sostener la estrategia anterior se introduce una caché de bloques descomprimidos. Su función es actuar como un conjunto de trabajo temporal
que almacena los bloques accedidos recientemente. Cuando una operación necesita un bloque, primero consulta la caché. Si el bloque está presente,
ocurre un acierto y se evita una nueva descompresión. Si el bloque no está presente, ocurre un fallo, el bloque se recupera desde almacenamiento
comprimido y se incorpora al conjunto de trabajo.

La política de reemplazo adoptada es del tipo \emph{Least Recently Used} (LRU). Bajo esta política, cuando la caché está llena y se necesita espacio
para un nuevo bloque, se expulsa el bloque cuyo último acceso sea el más antiguo. La elección de LRU es adecuada para este contexto porque muchas
secuencias de compuertas presentan localidad temporal: un bloque recién usado tiene mayor probabilidad de volver a ser requerido en el corto plazo
que uno accedido hace muchas operaciones.

Si el bloque expulsado fue modificado, debe recomprimirse antes de abandonar la caché. Si no fue modificado, puede descartarse sin costo de escritura.
Esta distinción entre bloques limpios y sucios evita recompresiones innecesarias y reduce la sobrecarga total del sistema.

\begin{figure}[H]
    \centering
    \fbox{\parbox{0.82\linewidth}{\centering
        Espacio reservado para un diagrama del funcionamiento de la caché LRU, mostrando aciertos, fallos, actualización de uso reciente y expulsión
        de bloques sucios o limpios.}}
    \caption{Esquema conceptual de la caché LRU de bloques descomprimidos}
    \label{fig:simulation-design-lru}
\end{figure}
